\chapter{Known limitations}
\label{ch:limitations}

This is the consolidated list.  Each entry states the limitation, the
measurement behind it where one exists, and the check a user can run.
The Tutorial has the same material in shorter form; the entries here
carry the numbers.

\section{Method-family limits}

\textbf{Contracted NEVPT2 accuracy.}  The perturbation layer inherits
the standard contracted-NEVPT2 approximations; the reference set's own
MAE for the method is the accuracy band claimed here, not a target
this package improves on.

\textbf{Valence--Rydberg mixing.}  $\pi\pi^*$ states near Rydberg
states move by up to an electron volt when diffuse functions enter
the basis.  Measured on ethylene at a fixed geometry: the singlet
$\pi\pi^*$ vertical excitation moves from $9.27$ to $8.09$~eV between
6-31G and 6-31+G* in the same three-orbital window, while the triplet
moves the opposite way ($4.83$ to $5.68$~eV); the singlet--triplet
gap, which the reference fixes near $3.4$~eV, collapses from $4.44$
to $2.41$~eV.  A gap collapse in the wrong direction is the signature
of a mixed-in state, not a converged one.  Check: re-run with and
without diffuse functions.

\textbf{Minimal-window limits for multiconfigurational
$\pi\pi^*$ states.}  The benchmark record shows a uniform positive
deviation of the $\pi\pi^*$ family ($+0.6$ to $+1.5$~eV) that
survives every window correction tried.  Check: widen the window by
one occupied orbital; a deviation that shrinks is window-limited.

\textbf{No spin--orbit coupling; scalar-relativistic with ECPs only.}
Fine structure is out of scope.
\textbf{Translated functionals beyond tPBE/tPBE0.}  The on-top kernel is a closed-form PBE, written out because the libxc wrapper cannot be traced; tPBE and its 0.25 hybrid tPBE0 are the covered pair.  Another translated functional needs its own closed-form kernel --- a one-line change in principle, a new validation in practice.  Below $\rho$ ~ 1e-13 the kernel differs from libxc internal guards by up to ~1e-4 absolute, where the grid weights put the integral contribution below 1e-25.


\section{Implementation limits}

\textbf{Dense diagonalization.}  The CI step materializes and
diagonalizes the full determinant space; the practical ceiling is a
few thousand determinants per root calculation.  A differentiable
iterative eigensolver is future work.  The (lambda, N) DVS scan inherits this ceiling: candidate windows are refused with their determinant count recorded once they exceed \texttt{max\_ndet} (default 2000), which places the published (12,12)/(14,14) windows beyond the in-stack CI by design --- those remain engine-side.

\textbf{Frozen on-top quadrature.}  Geometry derivatives of the on-top energy omit the grid-velocity term by construction; the calibration bounds it at 7.3e-5~Eh/bohr on CO/6-31G at grid level~3, falling to 2.5e-5 at level~4.


\textbf{Second-derivative cost.}  The implicit route requires the
orbital Hessian in memory and does not exploit sparsity; the unrolled
route scales with the iteration count.  Neither is competitive with
analytic second derivatives of single-reference methods.

\textbf{Memory and compile profiles of the differentiated chain.}
Intermediates are broadcast over the rotation-parameter count.  The
measured record that shaped the current design: an explicit Hessian
at about five hundred parameters requested a 402~GB allocation in a
single request; the Hessian-vector-product stages that replaced it
carry one tangent per pass, and their first evaluation still spends
minutes in XLA compilation at that size.  Check: watch resident
memory during the first optimization step, and leave headroom on
shared machines.

\subsection{Degenerate pairs at a stationary point}

The eigensolver's cluster tolerance must separate two scales: the
floating-point noise that splits symmetry-degenerate pairs at a
rotated stationary point (measured between $10^{-8}$ and
$1.2\times10^{-5}$ on the systems in this work) and the physical gaps
of the same spectrum ($1.9\times10^{-2}$ and above).  A tolerance
below the noise exposes the first and second derivatives to a bare
inverse gap; one above the physical gaps merges distinct states.  The
window measured here spans $10^{-5}$ to $10^{-2}$, and there is no
admissible value when a physical gap falls inside the noise band ---
the fallback for that case is the decomposition-free Lanczos kernel.
Check: if a linear or otherwise degenerate system returns unstable
derivatives, this tolerance is the first knob to inspect.

\section{Protocol limits}

\textbf{Multiple stationary points.}  The orbital energy is
non-convex; different starting points reach different stationary
points, and the selection is not stable even at fixed settings ---
water at 6-31G has two minima $13$~kcal/mol apart, and the
reference-orbital start reaches a saddle.  Report the stationarity
certificate and the stationary-point spread with every multireference
number.

\textbf{NEVPT2 convergence sensitivity.}  The correction depends on
the orbitals more sensitively than the energy; energy-based gates are
insufficient when the target is the correction.  Measured: orbital
stalls that move the correction by several kcal/mol across runs whose
energies agree to the microhartree.  Use the stability of the
correction across independent runs together with a tightened
orbital-gradient gate.

\textbf{Window rules and bookkeeping.}  The frontier audit covers the
patterns met so far; the electron count follows the window
($n_\mathrm{ele} = N_\mathrm{ele} - 2 n_\mathrm{core}$) and a
hard-coded count produces a wrong ground state that the energy alone
does not flag.
\textbf{One-step versus relaxed MC-PDFT.}  The one-step orbitals are not stationary for the on-top functional (measured |grad| ~ 3e-3 at the CASSCF point); a relaxed minimization lies 3.0e-4~Eh below on H2/6-31G.  The two numbers answer different questions --- papers using the one-step protocol must not be compared with relaxed values without saying which is which.

\textbf{On-top columns amplify orbital-point differences.}  The per-state on-top energy reacts to the state-averaged orbital point at first order, where the CASSCF column reacts at second order (slope ~3e-3 per radian at a CASSCF point).  Two runs whose state-averaged energies agree to 1e-7~Eh can differ by ~1e-5~eV in the on-top column; treat that as the column's reproducibility floor, not as noise.

